\documentclass[11pt]{article}
\usepackage{amsmath,amssymb,booktabs,geometry,graphicx,hyperref,url}
\usepackage{subcaption,placeins}
\hypersetup{hidelinks}
\geometry{margin=1in}
\setlength{\parskip}{0.45em}
\setlength{\parindent}{0pt}
\title{Stage 3: Boundary Extraction, Delta/Gamma, and Diagnostics}
\author{Free-Boundary-Aware Risk Surfaces for American Options}
\date{\today}
\begin{document}
\maketitle

\section{Stage 3: Boundary Extraction, Delta/Gamma, and Diagnostics}

\subsection{Purpose}
This stage extends the validated CN/PSOR benchmark with explicit
early-exercise boundary extraction, finite-difference Delta and Gamma surfaces,
strict Greek masks, and diagnostic tables for an American put and a
dividend-paying American call. Delta measures local price sensitivity to spot,
whereas Gamma measures the change in Delta and is therefore more sensitive to
grid noise and nonsmooth payoffs. This remains a classical-solver stage and
provides no neural-surrogate evidence.

\subsection{Method}
The experiment solves the base American put and dividend-paying American call
configurations on the default grid \(M=N=120\). For each option type, it
computes the continuation premium \(C=U-\Phi\), extracts the transition between
exercise and continuation, computes Delta and Gamma from the price surface, and
builds named masks around delicate regions.

\subsection{Mathematical Formulation}
The continuation premium is
\[
C(S,\tau)=U(S,\tau)-\Phi(S).
\]
For a put, the exercise region lies at low spot values, and the boundary is the
transition from \(C=0\) to \(C>0\). For a dividend-paying call, early exercise
may occur at high spot values, so the orientation is reversed. A missing
transition is recorded explicitly rather than converted into a fabricated
boundary value.

The Greeks are computed by finite differences:
\[
\Delta_i^n \approx \frac{U_{i+1}^n-U_{i-1}^n}{2\Delta S},\qquad
\Gamma_i^n \approx \frac{U_{i+1}^n-2U_i^n+U_{i-1}^n}{(\Delta S)^2}.
\]
Because Gamma is a second derivative, it amplifies interpolation error and the
nonsmoothness at the payoff kink and free boundary. The report therefore
separates all-grid summaries from strict-mask summaries; a Gamma spike in an
excluded band is a numerical warning, not automatically a solver failure.

\subsection{Implementation Details}
Boundary extraction now records a structured status for each time level:
\[
\texttt{boundary\_found},\quad \texttt{no\_exercise},\quad \texttt{all\_exercise}.
\]
The output CSV records boundary value, status, exercise node count, and continuation node count.

Greek masks are:
\begin{itemize}
  \item \texttt{near\_kink}: nodes near $S=K$;
  \item \texttt{near\_boundary}: nodes near the extracted free boundary;
  \item \texttt{near\_zero\_spot}: nodes too close to zero spot for stable relative interpretation;
  \item \texttt{strict}: nodes outside those delicate regions.
\end{itemize}

\subsection{Validation and Diagnostics}
For each option, the stage writes \path{boundary_curve.csv},
\path{greek_summary.csv}, \path{mask_summary.csv},
\path{diagnostics_summary.json}, and \path{surface_greeks.npz}. It also saves
price, Delta, Gamma, and boundary figures under
\path{results/figures/stage03_boundary_greeks/}.

\begin{table}[h]
\centering
\begin{tabular}{lrrrrr}
\toprule
Case & Obstacle & Complementarity & Monotonicity & Convexity & Strict mask\\
\midrule
Put & 0 & $4.0088\times10^{-12}$ & 0 & $2.22\times10^{-16}$ & 0.9669\\
Dividend call & 0 & $5.2696\times10^{-12}$ & $9.44\times10^{-14}$ & $4.50\times10^{-13}$ & 0.9670\\
\bottomrule
\end{tabular}
\caption{Base-case obstacle, complementarity, shape, and strict-mask diagnostics for the American put and dividend-paying call. Values are solver diagnostics on the \(120\times120\) grid, not surrogate results.}
\end{table}

\subsection{Results}
At the final time level, the American put boundary is $0.8K$ and the
dividend-paying American call boundary is $1.3333K$. Both have status
\texttt{boundary\_found}.

For the put, the all-grid Gamma maximum is $15.0$, while the strict-mask Gamma maximum is about $7.55$. Near the payoff kink the mean absolute Gamma is about $3.90$, and near the boundary it is about $1.63$. For the dividend-paying call, the all-grid Gamma maximum is also $15.0$, while the strict-mask Gamma maximum is about $7.50$. Near the kink the mean absolute Gamma is about $3.65$, and near the boundary it is about $0.45$.

The contrast between all-grid and strict-mask summaries shows that the largest
Gamma values are concentrated near mathematically delicate regions. The strict
mask therefore defines a more defensible region for later risk-surface
comparisons and supervised Delta labels without claiming universal Greek
stability.

% Auto-generated by experiments/19_report_visualization_enrichment.py
\section*{Visual Evidence Gallery}
\begin{figure}[!htbp]
\centering
\begin{subfigure}[t]{0.48\linewidth}
\centering
\includegraphics[width=\linewidth]{results/figures/report_enrichment/stage03_boundary_greeks_diagnostics_put_boundary_curve.pdf}
\caption{The put early-exercise boundary is extracted from low-spot exercise to continuation transitions in the CN/PSOR continuation premium.}
\end{subfigure}
\hfill
\begin{subfigure}[t]{0.48\linewidth}
\centering
\includegraphics[width=\linewidth]{results/figures/report_enrichment/stage03_boundary_greeks_diagnostics_call_boundary_curve.pdf}
\caption{The dividend-paying call boundary uses the opposite high-spot exercise orientation and is evaluated separately from put boundaries.}
\end{subfigure}
\end{figure}

\begin{figure}[!htbp]
\centering
\begin{subfigure}[t]{0.48\linewidth}
\centering
\includegraphics[width=\linewidth]{results/figures/report_enrichment/stage03_boundary_greeks_diagnostics_put_delta_surface.pdf}
\caption{Delta surfaces are interpreted together with masks around nonsmooth regions.}
\end{subfigure}
\hfill
\begin{subfigure}[t]{0.48\linewidth}
\centering
\includegraphics[width=\linewidth]{results/figures/report_enrichment/stage03_boundary_greeks_diagnostics_put_gamma_surface.pdf}
\caption{Gamma spikes near kinks and free boundaries are shown as diagnostic limitations rather than hidden failures.}
\end{subfigure}
\end{figure}

\FloatBarrier


\subsection{Failure Modes and Limitations}
Boundary extraction is still grid-resolution limited, so boundary values move in grid increments unless interpolation is supported by a clear premium transition. The current Greeks are finite-difference diagnostics from the solver surface; they are not smoothed Greeks and should not be over-interpreted near the payoff kink or free boundary. Rannacher smoothing is still not active in the solver, so near-maturity Gamma remains especially delicate.

\subsection{Conclusion and Next Step}
The Stage 3 gate passes for producing classical boundary and Greek diagnostics from the validated CN/PSOR solver. The next stage should quantify grid/domain sensitivity and runtime for these boundary and Greek quantities before any surrogate dataset construction.

\end{document}
